% pages-examples/tagged.tex   —— 配置②：带时间标签的条目
%
% 用于：课程、项目、经历这类「时间 + 内容」的清单。
%
% 用法要点：
%   - 用 \DeckTagged{标签}{内容}，标签会被撑到统一宽度，内容自然左对齐。
%   - 不要手打空格去凑对齐：标签字数不等（4 字 vs 3 字），
%     空格 padding 永远对不齐，全角半角混用会更乱。
%   - 标签宽度由 style.tex 的 \DeckTagWidth 控制（默认 5.6em，约 5 个全角字符）。
%   - 时间标签的词长尽量一致（都用"大二寒假"这种 4 字格式），读起来更整齐。

\DeckGap

\DeckTagged{大二寒假}{内容写在标签右侧，起点与其他行完全对齐}
\DeckTagged{大二暑假}{不需要手打空格凑宽度}
\DeckTagged{本学期}{三字标签也照样对齐}
